\documentclass[11pt,letterpaper]{article}

\usepackage[margin=1in]{geometry}
\usepackage{amsmath,amssymb,amsthm,bm}
\usepackage{booktabs}
\usepackage{enumitem}
\usepackage{xcolor}
\usepackage{graphicx}
\usepackage{tcolorbox}
\usepackage[colorlinks=true,linkcolor=blue!55!black,urlcolor=blue!55!black,citecolor=blue!55!black]{hyperref}

\graphicspath{{figs/}{notes/figs/}}

% Plain-language call-out box for readers new to quantum spin ice.
\newtcolorbox{intuition}[1][]{colback=blue!4,colframe=blue!35!black,
  boxrule=0.5pt,arc=2pt,left=6pt,right=6pt,top=4pt,bottom=4pt,
  fonttitle=\bfseries,title={In plain language},#1}

\newcommand{\Jzz}{J_{zz}}
\newcommand{\Jpm}{J_{\pm}}
\newcommand{\Heff}{H_{\mathrm{eff}}}
\newcommand{\Hclean}{H_{\mathrm{clean}}}
\newcommand{\gfour}{g_4}
\newcommand{\ghex}{g_{\mathrm{hex}}}
\newcommand{\bA}{\bm A}
\newcommand{\bd}{\bm d}
\newcommand{\bn}{\bm n}
\newcommand{\bN}{\bm N}
\newcommand{\bL}{\bm L}
\newcommand{\br}{\bm r}
\newcommand{\bw}{\bm w}
\newcommand{\bq}{\bm q}
\newcommand{\brho}{\bm \rho}
\newcommand{\bdelta}{\bm \delta}
\newcommand{\bDelta}{\bm \Delta}
\newcommand{\bphi}{\bm \phi}
\newcommand{\btheta}{\bm \theta}
\newcommand{\bW}{\bm W}
\newcommand{\Z}{\mathbb Z}
\newcommand{\Tr}{\operatorname{Tr}}
\newcommand{\sgn}{\operatorname{sgn}}
\newcommand{\ice}{\mathcal I}
\newcommand{\cycle}{\mathcal C}
\newcommand{\row}{\mathcal R}
\newcommand{\HB}{H_{\mathrm{band}}}
\newcommand{\Hmic}{H_{\mathrm{mic}}}

\theoremstyle{plain}
\newtheorem{proposition}{Proposition}
\newtheorem{lemma}{Lemma}
\newtheorem{corollary}{Corollary}
\theoremstyle{definition}
\newtheorem{definition}{Definition}
\newtheorem{remark}{Remark}

\title{Pedagogical notes on removing finite-size winding loops in exact
diagonalization of quantum spin ice}
\author{}
\date{\today}

\begin{document}
\maketitle

\begin{abstract}
Exact diagonalization (ED) is the natural numerical route to the $\pi$-flux
quantum spin ice (QSI) regime, where QMC has a sign problem.  The obstacle is
not the usual small-system limitation alone: periodic pyrochlore clusters
contain ice-preserving ring exchanges that close only by winding around the
torus.  The dominant artifact is a four-site winding loop generated at order
$\Jpm^2$, whereas the physical QSI photon ring exchange is the contractible
six-site hexagon generated at order $|\Jpm|^3$.  These notes derive the
transported-dipole label
\[
        \bdelta=\brho+\bN\bL,
\]
show why ordinary boundary-twist averaging does not remove the artifact, and
prove two exact structure theorems that collapse the character-average
protocol into its final form: at the minimal grid the character deformation is
unitarily equivalent to the bare Hamiltonian through the diagonal
polarization unitary $e^{2i\btheta\cdot\hat{\bm X}}$, so the entire
eight-point character average of the downfolded ice-band operator equals a
single untwisted downfolding followed by a \emph{polarization-parity mask}.
Homotopy fixes the scope: the torus is a $K(\Z^3,1)$, the winding class is
the complete topological invariant of a closed process, and the mask is the
maximal topologically defined cleaning.  The production protocol at any
$\Jpm$ is then \emph{graded Brillouin--Wigner downfolding}: the exact
ice-block resolvent operator $F(z)=PHP+PHQ(z-QHQ)^{-1}QHP$, masked entry by
entry, solved self-consistently.  It requires no eigenvectors, no band
selection, and no Gram matrix; its diagnostics remain meaningful at every
coupling, and its breakdown modes are physical (dissolution of the gauge
sector into the spinon continuum), not numerical.  Every claim is verified
against exact microscopic ED on the $16$-site cluster.
\end{abstract}

\tableofcontents

\section{The idea in one page (for readers new to spin ice)}

Before any equations, here is the whole story in words and pictures.

\begin{itemize}
\item \textbf{What the material is.}  ``Quantum spin ice'' (QSI) is a magnet
in which the atomic magnetic moments sit on a lattice of corner-sharing
tetrahedra (the \emph{pyrochlore} lattice).  On every tetrahedron the moments
obey a local rule: two point in and two point out
(Fig.~\ref{fig:ice_rule}).  This is the same counting rule that governs the
proton positions in ordinary water ice, hence the name.  There are
astronomically many ways to satisfy it, so the system has a huge
ground-state degeneracy.

\item \textbf{What makes it ``quantum''.}  A small quantum-mechanical term
lets pairs of neighbouring moments flip.  Acting over and over, these flips
let the system tunnel between different two-in--two-out patterns.  The
smallest tunnelling move that respects the ice rule everywhere is a
\emph{six-spin ring exchange} around a hexagon
(Fig.~\ref{fig:diamond_hexagon} and Fig.~\ref{fig:ring_exchange}).  Remarkably,
the collective dynamics of these ring flips is described by an emergent
\emph{light}: QSI hosts its own artificial photon.  This is what makes the
material exciting, and it is what we want to compute.

\item \textbf{Why the computation is hard.}  We solve the model on a computer
by exact diagonalization (ED), which requires a \emph{finite} block of the
lattice with periodic boundaries---mathematically a doughnut, or
\emph{torus}.  On a torus, some sequences of flips close up not because they
form a real loop in space, but only because they run off one face and come
back on the opposite face (Fig.~\ref{fig:torus_winding}).  The shortest such
``wrap-around'' move involves only \emph{four} spins, and---crucially---it
appears at a lower order in perturbation theory than the physical six-spin
hexagon.  It is a pure artifact of the finite box, yet it is numerically
\emph{larger} than the physics we are after (Fig.~\ref{fig:scale_sep}).

\item \textbf{The symptom.}  Left untreated, this artifact puts the
low-temperature specific-heat peak---one of the cleanest experimental
fingerprints of the emergent photon---in the wrong place, at the artifact
scale $g_4\propto J_\pm^2$ instead of the photon scale
$g_{\mathrm{hex}}\propto|J_\pm|^3$ (Figs.~\ref{fig:ct_twopeak}
and~\ref{fig:specific_heat_all_cases}).

\item \textbf{The cure.}  The physical hexagon closes in real space; the
artifact only closes through the wall of the box.  We give a single
gauge-invariant number, the \emph{transported dipole}
$\bdelta=\brho+\bN\bL$, that is exactly zero for real-space processes and
nonzero for wrap-around ones (Fig.~\ref{fig:real_loops}).  In perturbation
theory this says to keep only the zero-transport rows.  Nonperturbatively,
two exact theorems reduce the corresponding operation to something almost
embarrassingly simple: express the downfolded low-energy operator in the ice
basis and \emph{delete every matrix element whose polarization difference is
odd} (the parity mask).  At any coupling, the object to mask is the exact
Brillouin--Wigner ice-block operator $F(z)$, which exists whether or not the
low band is still ice-like.  We also show why simpler fixes, such as
boundary-twist averaging or averaging thermodynamic curves, do not remove
the artifact (Fig.~\ref{fig:twist_fail}) --- indeed for thermodynamic curves
we prove that they \emph{cannot}.
\end{itemize}

The rest of these notes makes each of these statements precise and checks
every number with an independent, from-scratch script.

\section{The physical question}

The QSI Hamiltonian considered here is the nearest-neighbor XXZ model in
local pyrochlore frames,
\begin{equation}
H =
\Jzz\sum_{\langle ij\rangle} S_i^zS_j^z
-\Jpm\sum_{\langle ij\rangle}
\left(S_i^+S_j^-+S_i^-S_j^+\right),
\qquad |\Jpm|\ll \Jzz,\quad \Jzz>0 .
\label{eq:model}
\end{equation}
The Ising part has an extensively degenerate ground manifold: every
tetrahedron has two spins pointing in and two spins pointing out.  This is
the ice manifold $\ice$.

For $\Jpm>0$, the off-diagonal matrix elements in the $S^z$ basis are
non-positive, and standard QMC can be sign-problem-free.  For $\Jpm<0$,
the off-diagonal matrix elements have the opposite sign.  This is the
frustrated $\pi$-flux QSI sector.  It is precisely the sector where ED is
most valuable, because QMC no longer gives an unbiased finite-temperature
benchmark.

The central difficulty is that ED necessarily uses a finite periodic
cluster.  Such a cluster is a torus.  On a torus, paths can close by
periodic identification even when they do not close in real space.  Some of
these winding paths preserve the ice rule and therefore appear directly in
the low-energy gauge Hamiltonian.  The shortest such processes on the
standard 16-site cubic and 32-site FCC clusters are four-site loops.  They
are not physical bulk plaquettes.  Nevertheless, they strongly affect the
low-temperature specific heat and the low-frequency $S^{zz}(\bq,\omega)$.

\section{Ice variables and the energy cost of defects}

\begin{figure}[htbp]
\centering
\includegraphics[width=0.95\textwidth]{fig_ice_rule.pdf}
\caption{\textbf{The ice rule.}  Each spin lives on a corner of a
tetrahedron and points either into the tetrahedron (red) or out of it
(blue).  \emph{Left:} the low-energy ``ice'' configuration has two in and two
out, so the net Ising charge $Q_t=\sum_{i\in t}S_i^z$ vanishes and the
tetrahedron pays no energy.  \emph{Middle, right:} tipping one spin makes the
tetrahedron three-in--one-out or one-in--three-out, with $|Q_t|=1$; by
Eq.~\eqref{eq:tet_energy} this costs $\Jzz/2$.  These charged tetrahedra are
the ``spinons'' of quantum spin ice.}
\label{fig:ice_rule}
\end{figure}

For a tetrahedron $t$, define the Ising charge
\begin{equation}
Q_t=\sum_{i\in t} S_i^z .
\end{equation}
For spin $1/2$, $S_i^z=\pm 1/2$.  The Ising energy on one tetrahedron is
\begin{equation}
\sum_{\langle ij\rangle\in t} S_i^zS_j^z
=\frac12\left[\left(\sum_{i\in t}S_i^z\right)^2-\sum_{i\in t}(S_i^z)^2\right]
=\frac12 Q_t^2-\frac12 .
\label{eq:tet_energy}
\end{equation}
Thus, up to a constant,
\begin{equation}
H_0=\frac{\Jzz}{2}\sum_t Q_t^2 .
\label{eq:charge_h0}
\end{equation}
The ice rule is $Q_t=0$ on every tetrahedron.  A $3$-in--$1$-out or
$1$-in--$3$-out tetrahedron has $|Q_t|=1$ and costs $\Jzz/2$ relative to an
ice tetrahedron.

Now act with a transverse exchange
\[
S_i^+S_j^-+S_i^-S_j^+
\]
on an antiparallel nearest-neighbor pair.  The two sites $i$ and $j$ lie on
one common tetrahedron, where exchanging an up and a down spin preserves
$Q_t=0$.  Each site also belongs to one other tetrahedron.  Those two
neighboring tetrahedra become charged, each costing $\Jzz/2$.  Therefore a
single nearest-neighbor exchange out of the ice manifold creates two
defects with total cost
\begin{equation}
\Delta_{\mathrm{1hop}}=\Jzz .
\label{eq:one_hop_cost}
\end{equation}
This simple denominator is the origin of the perturbative scales below.

\begin{intuition}
A single spin flip is like pushing one domino: it necessarily unbalances
the two tetrahedra that share that spin, creating a pair of oppositely
charged defects that together cost $\Jzz$.  To get back into the ice manifold
you must flip more spins so that every tetrahedron is balanced again.  The
cheapest way to do that around a closed path is what the next sections count.
\end{intuition}

\section{Degenerate perturbation theory in the ice manifold}

Let $P$ be the projector onto the ice manifold and $Q=1-P$.  Write
\[
H=H_0+V,\qquad
V=-\Jpm\sum_{\langle ij\rangle}
\left(S_i^+S_j^-+S_i^-S_j^+\right).
\]
The ice energy is shifted to zero.  The non-ice resolvent is
\begin{equation}
R=Q\frac{1}{E_{\ice}-H_0}Q=-QH_0^{-1}Q .
\end{equation}
Because a single exchange leaves the ice manifold, $PVP=0$.  The first
nonzero terms are therefore
\begin{align}
\Heff^{(2)} &= PVRVP, \label{eq:heff2}\\
\Heff^{(3)} &= PVRVRVP. \label{eq:heff3}
\end{align}
For a path $s\to m_1\to \cdots \to t$, its contribution has the form
\begin{equation}
\langle t|\Heff^{(k)}|s\rangle
=
\sum_{\mathrm{paths}}
\frac{\langle t|V|m_{k-1}\rangle\cdots
\langle m_1|V|s\rangle}
{(E_\ice-E_{m_1})\cdots(E_\ice-E_{m_{k-1}})} .
\label{eq:path_pt}
\end{equation}
Each matrix element of $V$ contributes a factor $-\Jpm$, and each
intermediate energy denominator is negative.

\begin{remark}
The standalone recomputation stores each perturbative process as a row
\[
        r=(s,t,\bN,c,k),
\]
meaning that the row connects ice states $s$ and $t$, appears at order $k$,
has coefficient $c$ when $\Jpm=1$, and has accumulated periodic image
transport $\bN$.  At physical coupling, the row contributes
$c\Jpm^k|t\rangle\langle s|$.
\end{remark}

\section{Why the physical QSI ring exchange is a hexagon}

\begin{figure}[htbp]
\centering
\includegraphics[width=0.92\textwidth]{fig_diamond_hexagon.pdf}
\caption{\textbf{The smallest real ring exchange has six spins.}  Each
tetrahedron is a vertex of the diamond lattice (circles); each spin is a link
of that lattice (squares).  A ring exchange must flip an alternating chain of
spins around a closed loop of tetrahedra.  \emph{Left:} the diamond lattice is
bipartite (two colours, links only join different colours), so its shortest
loop uses six links---a six-spin hexagon.  \emph{Right:} a four-spin ring
would require a four-loop on the diamond lattice, which would need two
same-colour vertices to be adjacent; no such loop exists in the bulk.  This is
the geometric reason a contractible four-site ice loop is impossible.}
\label{fig:diamond_hexagon}
\end{figure}

The pyrochlore lattice is the line graph of the diamond lattice: a
pyrochlore site corresponds to a diamond bond, and a pyrochlore
tetrahedron corresponds to a diamond vertex.  An ice configuration assigns
a divergence-free electric field to diamond bonds.  A ring exchange in the
ice manifold must flip an alternating set of spins around a closed path so
that the divergence at every touched diamond vertex remains zero.

\begin{lemma}[Shortest contractible ice loop]
The shortest nontrivial contractible ice-preserving ring exchange on the
infinite pyrochlore lattice is a six-site hexagon.
\end{lemma}

\begin{proof}
An ice-preserving loop must enter and leave every touched tetrahedron, so
on the diamond lattice it corresponds to a closed non-backtracking loop.
The diamond lattice is bipartite and has no triangles.  It also has no
four-cycles; its elementary loops are length six.  Therefore the shortest
contractible closed diamond loop has length six.  Mapping back to the
pyrochlore line graph gives a six-site pyrochlore hexagon.  A triangle of
the pyrochlore graph lies inside one tetrahedron and is not an alternating
ice ring.  A contractible four-site ice ring would map to a diamond
four-cycle, which does not exist.
\end{proof}

\subsection{Third-order hexagon amplitude}

\begin{figure}[htbp]
\centering
\includegraphics[width=0.95\textwidth]{fig_ring_exchange.pdf}
\caption{\textbf{The physical ring exchange (the emergent-photon move).}
Starting from an alternating up/down pattern on a hexagon, flipping all three
alternating pairs turns each up into a down and vice versa, returning to the
ice manifold.  Because the intermediate steps leave the ice manifold, the
amplitude is third order in $J_\pm$, with magnitude
$g_{\mathrm{hex}}=12|J_\pm|^3/J_{zz}^2$.  The coherent superposition of these
ring flips across the lattice is the artificial photon of quantum spin ice.}
\label{fig:ring_exchange}
\end{figure}

Consider a flippable hexagon with alternating spins.  To move from one ice
configuration to the flipped configuration, one must exchange three
alternating nearest-neighbor pairs.  After the first exchange the system
has two charged tetrahedra and costs $\Jzz$.  After the second exchange it
still has a pair of charged tetrahedra and costs $\Jzz$.  Thus each
ordering gives a denominator $\Jzz^2$.

For one orientation there are $3!$ orderings.  Including the Hermitian
conjugate orientation gives the conventional ring term
\begin{equation}
\Heff^{\mathrm{hex}}
=-\ghex\sum_{\mathrm{hex}}
\left(
S_1^+S_2^-S_3^+S_4^-S_5^+S_6^-+\mathrm{h.c.}
\right),
\end{equation}
with magnitude
\begin{equation}
\boxed{
\ghex=\frac{12|\Jpm|^3}{\Jzz^2}.
}
\label{eq:ghex}
\end{equation}
The sign of the effective ring term follows $\sgn(\Jpm^3)$.  This sign is
what distinguishes the $0$-flux and $\pi$-flux gauge sectors.

\section{Why a four-site loop is a finite-size artifact}

\begin{figure}[htbp]
\centering
\includegraphics[width=0.9\textwidth]{fig_torus_winding.pdf}
\caption{\textbf{A periodic cluster is a torus.}  Opposite walls of the
simulation box are identified (glued), so the box is topologically a
doughnut.  \emph{Left:} a loop that closes inside the box has winding
$\bw=0$; it is a genuine real-space process such as the physical hexagon.
\emph{Right:} a path can also close by leaving one wall and re-entering the
opposite wall.  Its winding $\bw\neq0$; it exists \emph{only} because of the
periodic gluing and has no counterpart in the infinite lattice.}
\label{fig:torus_winding}
\end{figure}

A periodic cluster is not $\mathbb R^3$; it is the flat torus
\begin{equation}
T^3=\mathbb R^3/(\Z^3\bL) ,
\end{equation}
where we use row vectors and the rows of $\bL$ are the three cluster
translation vectors.  Thus $\bm m\bL$, with $\bm m\in\Z^3$, is a periodic
translation.  For the cubic $L_1\times L_2\times L_3$ cluster,
\[
\bL=\begin{pmatrix}
L_1&0&0\\
0&L_2&0\\
0&0&L_3
\end{pmatrix}.
\]
For the FCC primitive cluster, the rows are
$L_1(\tfrac12,\tfrac12,0)$, $L_2(\tfrac12,0,\tfrac12)$, and
$L_3(0,\tfrac12,\tfrac12)$.  All coordinates in this note are measured in
conventional cubic-cell units.  The pyrochlore basis sites lie on the
quarter-lattice,
\[
\br_i\in \frac14\Z^3,
\]
which is the reason factors of $4$ appear below.

A bond
from site $i$ to site $j$ is represented by a minimum-image displacement
\begin{equation}
\bd_{ij}=\br_j-\br_i-\bn_{ij}\bL,
\qquad \bn_{ij}\in\Z^3 .
\end{equation}
For a closed path $\cycle$ on the periodic graph, define its winding vector
\begin{equation}
\bw_\cycle=\sum_{(ij)\in\cycle}\bn_{ij}.
\label{eq:winding}
\end{equation}
If $\bw_\cycle=0$, the lifted path closes in real space and is
contractible.  If $\bw_\cycle\ne0$, the path closes only because opposite
faces of the cluster have been identified.

\begin{proposition}[No physical four-loop]
Every ice-preserving four-cycle on the standard $1\times1\times1$ cubic
and $2\times2\times2$ FCC pyrochlore clusters is a winding torus loop, not
a contractible bulk process.
\end{proposition}

\begin{proof}
The geometric lemma above rules out a contractible four-site ice loop on
the infinite pyrochlore lattice.  Therefore any four-site ice loop observed
on a periodic cluster must use the periodic identification.  Equivalently,
its lift to $\mathbb R^3$ fails to close and has nonzero winding
$\bw_\cycle$.  The explicit independent enumeration below verifies this
for the two clusters used in the ED calculation.
\end{proof}

\subsection{Second-order four-loop amplitude}

\begin{figure}[htbp]
\centering
\includegraphics[width=0.95\textwidth]{fig_pt_ladder.pdf}
\caption{\textbf{Counting the powers of $J_\pm$.}  Every transverse exchange
adds one factor of $J_\pm$ and pushes the system one virtual step above the
ice manifold, where it costs $\sim\Jzz$.  \emph{Left:} the winding four-loop
returns to the ice manifold after just two exchanges (one intermediate
tetrahedron pair), giving $g_4=4J_\pm^2/\Jzz$.  \emph{Right:} the hexagon
needs three exchanges (two intermediate steps), giving the weaker,
higher-order $g_{\mathrm{hex}}=12|J_\pm|^3/\Jzz^2$.  The extra rung is the
entire reason the artifact is parametrically larger at small $J_\pm$.}
\label{fig:pt_ladder}
\end{figure}

A winding four-loop has four alternating sites and can be flipped by two
nearest-neighbor exchanges.  The two exchanged bonds form a perfect
matching of the four sites.  There are two matchings, and for each matching
there are two time orderings.  Each path has one intermediate state with
energy cost $\Jzz$.  Thus the magnitude of the second-order matrix element
is
\begin{equation}
4\frac{\Jpm^2}{\Jzz}.
\end{equation}
In the ring-exchange convention used here,
\begin{equation}
\boxed{
\gfour=\frac{4\Jpm^2}{\Jzz}.
}
\label{eq:g4}
\end{equation}
The sign is convention dependent, but the scale is not.

The scale separation is
\begin{equation}
\boxed{
\frac{\gfour}{\ghex}
=
\frac{4\Jpm^2/\Jzz}{12|\Jpm|^3/\Jzz^2}
=
\frac{\Jzz}{3|\Jpm|}.
}
\label{eq:scale_separation}
\end{equation}
Thus the artifact becomes more severe precisely in the perturbative regime
where the QSI description is best controlled.

\begin{figure}[htbp]
\centering
\includegraphics[width=0.62\textwidth]{fig_scale_sep.pdf}
\caption{\textbf{The artifact dominates exactly where we want to work.}  The
finite-size four-loop scale $g_4$ (orange) and the physical photon scale
$g_{\mathrm{hex}}$ (purple) both vanish as $J_\pm\to0$, but their ratio
$g_4/g_{\mathrm{hex}}=\Jzz/(3|J_\pm|)$ \emph{grows}.  At the representative
value $|J_\pm|=0.05\,\Jzz$ the artifact is already $\sim7\times$ larger than
the physics it would otherwise swamp.}
\label{fig:scale_sep}
\end{figure}

\section{Independent numerical verification}

The script
\[
\texttt{notes/recompute\_finite\_size\_artifact.py}
\]
does not import the old \texttt{twist\_qsi\_demo/scripts} code and does not
inspect git history.  It independently:
\begin{enumerate}[label=(\roman*)]
\item builds the pyrochlore positions, nearest-neighbor graph, tetrahedra,
and periodic image vectors;
\item enumerates all two-in--two-out ice states;
\item enumerates ice-preserving four-cycles and hexagons;
\item constructs $H_2=PVRVP$ and $H_3=PVRVRVP$ from path sums;
\item applies either no projector or the zero-transport projector derived
below;
\item diagonalizes the resulting ice-manifold Hamiltonian and computes
$C(T)$.
\end{enumerate}
The raw output is saved in
\texttt{notes/recomputed\_finite\_size\_artifact.json}.

\subsection{Geometry census}

\begin{center}
\begin{tabular}{lrrrrr}
\toprule
cluster & sites & ice states & 4-cycles & contractible 4-cycles
& contractible hexagons \\
\midrule
cubic $1\times1\times1$ & 16 & 90 & 36 & 0 & 16 \\
FCC $2\times2\times2$ & 32 & 2970 & 24 & 0 & 32 \\
\bottomrule
\end{tabular}
\end{center}

The same enumeration finds 48 wrapping hexagons on the 16-site cluster and
96 wrapping hexagons on the 32-site cluster.  Those are also finite-size
torus processes, but they appear at the same order as the physical
hexagons.  The four-loops are the dominant contamination because they
appear one order earlier.

\section{Why the low-temperature specific heat peak is displaced}

For an exact spectrum $\{E_n\}$, shifted so $E_0=0$,
\begin{equation}
Z(T)=\sum_n e^{-E_n/T},\qquad
\langle E^p\rangle_T=\frac{1}{Z}\sum_n E_n^p e^{-E_n/T},
\end{equation}
and
\begin{equation}
C(T)=\frac{\langle E^2\rangle_T-\langle E\rangle_T^2}{T^2}.
\label{eq:specific_heat}
\end{equation}
For a two-level system with splitting $\Delta$,
\begin{equation}
C_{\mathrm{2lev}}(T)
=\frac{(\Delta/T)^2 e^{-\Delta/T}}
{\left(1+e^{-\Delta/T}\right)^2},
\end{equation}
whose maximum is at $T=O(\Delta)$.  A many-level finite cluster has a more
complicated spectrum, but the lower peak still tracks the dominant
low-energy splitting scale in the ice manifold.

Therefore:
\begin{align}
T_{\mathrm{peak}}^{\mathrm{bare}} &\sim \gfour\propto \Jpm^2,
\qquad\text{if winding four-loops are present},\\
T_{\mathrm{peak}}^{\mathrm{clean}} &\sim \ghex\propto |\Jpm|^3,
\qquad\text{after winding transport is removed}.
\end{align}

\begin{figure}[htbp]
\centering
\includegraphics[width=0.66\textwidth]{fig_ct_twopeak.pdf}
\caption{\textbf{The low-temperature specific-heat peak is displaced by the
artifact.}  Exact specific heat of the $16$-site cubic ice-manifold
Hamiltonian at $J_\pm=-0.05\,\Jzz$.  The bare cluster (orange) peaks at the
artifact scale $g_4$; after the $\bdelta=0$ projection (green) the peak drops
onto the physical photon scale $g_{\mathrm{hex}}$---nearly an order of
magnitude lower in temperature.  Since the specific heat is a nonlinear
function of the spectrum, this shift cannot be undone by averaging $C(T)$
after the fact (Sec.~\ref{sec:qed-spine}).}
\label{fig:ct_twopeak}
\end{figure}

\begin{figure}[htbp]
\centering
\includegraphics[width=\textwidth]{fig_specific_heat_all_cases.pdf}
\caption{\textbf{Specific heat as a function of temperature for all
independently recomputed scenarios.}  Each panel shows the exact
ice-manifold specific heat per site, $C(T)/N$, from the standalone
Schrieffer--Wolff calculation.  The columns are the three computed
couplings, and the rows are the two ED clusters.  Orange is the bare
periodic effective Hamiltonian; green is the zero-transport
$(\bdelta=0)$ Hamiltonian.  Dotted orange/green vertical lines mark the
actual peak positions of the two curves, the black dash-dot line marks
$g_4$, and the purple dashed line marks $g_{\mathrm{hex}}$.  The full
curves show the same conclusion as the peak table: removing transported
winding processes shifts the low-temperature feature from the finite-size
scale toward the photon scale.}
\label{fig:specific_heat_all_cases}
\end{figure}

The independent calculation gives the following representative peaks:
\begin{center}
\begin{tabular}{llrrr}
\toprule
cluster & Hamiltonian & $\Jpm$ & $T_{\mathrm{peak}}$ &
scale diagnostic \\
\midrule
cubic $N=16$ & bare & $-0.05$ & 0.010056 &
$T_{\mathrm{peak}}/g_4=1.01$ \\
cubic $N=16$ & clean & $-0.05$ & 0.001604 &
$T_{\mathrm{peak}}/g_{\mathrm{hex}}=1.07$ \\
FCC $N=32$ & bare & $-0.05$ & 0.004180 &
$T_{\mathrm{peak}}/g_{\mathrm{hex}}=2.79$ \\
FCC $N=32$ & clean & $-0.05$ & 0.000782 &
$T_{\mathrm{peak}}/g_{\mathrm{hex}}=0.52$ \\
FCC $N=32$ & bare & $-0.10$ & 0.018053 &
$T_{\mathrm{peak}}/g_{\mathrm{hex}}=1.50$ \\
FCC $N=32$ & clean & $-0.10$ & 0.006252 &
$T_{\mathrm{peak}}/g_{\mathrm{hex}}=0.52$ \\
\bottomrule
\end{tabular}
\end{center}
For the 16-site cluster at $\Jpm=-0.05$, $g_4=4(0.05)^2=0.010$ and
$g_{\mathrm{hex}}=12(0.05)^3=0.0015$, so the scale identification is
especially transparent.

\begin{figure}[htbp]
\centering
\includegraphics[width=\textwidth]{fig_all_scenarios.pdf}
\caption{\textbf{All independently recomputed finite-size scenarios in one
place.}  \emph{(a)} The two ED clusters both contain winding four-loops
and wrapping hexagons in addition to the physical contractible hexagons.
\emph{(b)} Boundary-twist averaging keeps half of each winding channel on
both clusters, while the zero-transport projector removes them and keeps
the physical hexagon.  \emph{(c)} The low-temperature peak positions for
all computed clusters, flux signs, and projection choices.  Black ticks
mark $g_4$ and purple crosses mark $g_{\mathrm{hex}}$.  \emph{(d)} The same
peak data in photon units.  The cleaned spectra stay at an order-one
multiple of $g_{\mathrm{hex}}$, while the bare spectra are pushed upward by
the winding-loop contamination.}
\label{fig:all_scenarios}
\end{figure}

\section{Why ordinary boundary-twist averaging fails}

Boundary twists are often used to reduce finite-size effects.  Put phases
only on boundary-crossing bonds.  If a virtual path accumulates image
transport $\bN$, its row coefficient becomes
\begin{equation}
c_r(\bphi)=c_r e^{-i\bN_r\cdot\bphi}.
\end{equation}
The eight-corner average over $\phi_a\in\{0,\pi\}$ gives
\begin{equation}
\frac18\sum_{\phi_a=0,\pi}e^{-i\bN\cdot\bphi}
=\prod_{a=1}^3\frac{1+(-1)^{N_a}}{2}.
\label{eq:corner_kernel}
\end{equation}
This keeps rows with even $\bN$.

The subtle point is that the geometric ring exchange is not a single path.
It is a sum over virtual paths.  A winding four-loop has multiple
matchings and time orderings.  Some paths cross the boundary, while other
paths implementing the same geometric flip can have $\bN=0$ in the
boundary gauge.  The corner average therefore kills only part of the
coefficient.

\begin{figure}[htbp]
\centering
\includegraphics[width=0.66\textwidth]{fig_twist_fail.pdf}
\caption{\textbf{Boundary-twist averaging is not enough.}  Fraction of each
geometric channel's coupling that survives on the $16$-site cubic cluster.
Ordinary boundary-twist (corner) averaging (grey) leaves \emph{half} of every
winding artifact---because the ring exchange is a sum over virtual paths and
some of those paths do not cross the boundary in the chosen gauge.  The
gauge-invariant $\bdelta=0$ projection (green) removes the winding four-loop
and wrapping hexagon completely while keeping the physical hexagon at full
strength.}
\label{fig:twist_fail}
\end{figure}

The standalone recomputation classifies the off-diagonal rows by geometric
channel and gives:
\begin{center}
\begin{tabular}{lccc}
\toprule
channel & corner average kept & $N=0$ kept & zero-transport kept \\
\midrule
4-loop, cubic $N=16$ & 0.5 & 0.5 & 0 \\
wrapping hexagon, cubic $N=16$ & 0.5 & 0.5 & 0 \\
contractible hexagon, cubic $N=16$ & 1 & 1 & 1 \\
4-loop, FCC $N=32$ & 0.5 & 0.5 & 0 \\
wrapping hexagon, FCC $N=32$ & 0.5 & 0.5 & 0 \\
contractible hexagon, FCC $N=32$ & 1 & 1 & 1 \\
\bottomrule
\end{tabular}
\end{center}
Thus ordinary boundary-twist averaging reduces the winding artifact but
does not remove it.

\section{The transported dipole and the zero-transport theorem}
\label{sec:transported-dipole}

The correct object is not the boundary-crossing count of one chosen gauge.
It is the gauge-invariant net displacement transported by the effective
operator.

\begin{definition}[Home-cell dipole]
For a perturbative row connecting ice states $s$ and $t$, define
\[
R_+(t,s)=\{i:S_i^z(t)-S_i^z(s)>0\},\qquad
R_-(t,s)=\{j:S_j^z(t)-S_j^z(s)<0\}.
\]
The set $R_+$ contains spins raised by the effective process; $R_-$ contains
spins lowered by the effective process.  The home-cell dipole is
\begin{equation}
\brho_{t,s}=\sum_{i\in R_+(t,s)}\br_i-\sum_{j\in R_-(t,s)}\br_j .
\label{eq:rho_def}
\end{equation}
\end{definition}

\begin{definition}[Transported dipole]
Let $\bN_r\in\Z^3$ be the accumulated image vector of the virtual path.  In row
vector convention, $\bN_r\bL$ is the net periodic translation needed when the
path is lifted from the torus back to the covering space $\mathbb R^3$.  The
transported dipole is
\begin{equation}
\boxed{
\bdelta_r=\brho_{t,s}+\bN_r\bL .
}
\label{eq:delta_def}
\end{equation}
\end{definition}

The formula is easiest to understand for one microscopic hop.  The operator
$S_i^+S_j^-$ raises site $i$ and lowers site $j$, so its home-cell dipole is
$\br_i-\br_j$.  If the nearest-neighbor bond was represented using the periodic
image vector $\bn_{ij}$, then the lifted displacement of the spin transfer is
\[
\br_i-\br_j+\bn_{ij}\bL=-\bd_{ij}.
\]
The many-step perturbative row adds these lifted increments.  Thus
$\bdelta_r$ is not a new dynamical charge.  It is the net real-space
displacement carried by the virtual spin transfers after the periodic images
have been undone.

For the order-$\Jpm^2$ and order-$\Jpm^3$ ice-manifold rows explicitly
enumerated below, the possible transported dipoles are coarser than the
microscopic quarter-lattice: their components lie in $\frac12\Z$.  This is the
only set of virtual processes that controls the unwanted four-loop peak and
the physical hexagon peak in the present campaign.

\begin{figure}[htbp]
\centering
\includegraphics[width=0.95\textwidth]{fig_real_loops.pdf}
\caption{\textbf{The transported dipole tells the two loops apart, on the
real $16$-site cubic cluster.}  Both loops flip alternating spins (red raised,
blue lowered) and both return the system to the ice manifold.  \emph{Left:}
the contractible hexagon is balanced---the raised and lowered spins average to
the same point---so its transported dipole vanishes,
$\bdelta=\brho+\bN\bL=0$.  \emph{Right:} the winding four-loop must cross the
torus wall (dashed bonds); its raised and lowered spins are offset by half a
unit cell, giving $\bdelta=(0,0,-\tfrac12)\neq0$.  Selecting $\bdelta=0$ keeps
the left process and discards the right one.}
\label{fig:real_loops}
\end{figure}

\begin{lemma}[Gauge invariance of $\bdelta$]
The vector $\bdelta_r$ is independent of the choice of home cell and of the
gauge used to place the twist phases.
\end{lemma}

\begin{proof}
Move a site coordinate by a cluster translation, $\br_i\mapsto
\br_i+\bm m\bL$ with $\bm m\in\Z^3$.  If that site is raised, $\brho$
changes by $+\bm m\bL$; if it is lowered, $\brho$ changes by
$-\bm m\bL$.  The image bookkeeping of the virtual path changes by the
opposite integer vector, because the same physical hop is now represented
using a different periodic image.  Therefore $\bN\bL$ changes by the
negative of the change in $\brho$, and the sum $\brho+\bN\bL$ is unchanged.
\end{proof}

\begin{proposition}[Transported-dipole character projection]
For the order-$\Jpm^2$ and order-$\Jpm^3$ ice-manifold rows on the clusters
considered here, every transported dipole has components in $\frac12\Z$.  We
therefore define the integer row label
\begin{equation}
\boxed{\bDelta_r\equiv2\bdelta_r\in\Z^3.}
\label{eq:Delta_integer_label}
\end{equation}
This is only a change of units for these rows.  For example, the winding four-loop in
Fig.~\ref{fig:real_loops} has
$\bdelta=(0,0,-\tfrac12)$, hence $\bDelta=(0,0,-1)$.  The condition
$\bDelta=0$ is exactly the same condition as $\bdelta=0$.

Couple an auxiliary character field $\btheta$ to this integer transported
dipole.  A perturbative row $r$ acquires the phase
\begin{equation}
c_r(\btheta)=c_r e^{-i\btheta\cdot\bDelta_r}.
\label{eq:delta_character_phase}
\end{equation}
The operator-level average over $\btheta\in[0,2\pi)^3$ keeps exactly the rows
with $\bdelta_r=0$.
\end{proposition}

\begin{proof}
The characters of the abelian group $\Z^3$ are orthogonal:
\[
\overline{c_r}
=c_r\int_{[0,2\pi)^3}\frac{d^3\theta}{(2\pi)^3}
e^{-i\btheta\cdot\bDelta_r}
=c_r\,\delta_{\bDelta_r,0}.
\]
Since $\bDelta_r=2\bdelta_r$, the condition $\bDelta_r=0$ is equivalent to
$\bdelta_r=0$.
\end{proof}

\begin{remark}[How many discrete character points are needed?]
The exact integral over $\btheta$ can be replaced by a finite Fourier sum if
the grid separates the transported-dipole sectors present in the calculation.
For one component,
\[
\frac1M\sum_{m=0}^{M-1}
e^{-i(2\pi m/M)\Delta}
=
\begin{cases}
1,& \Delta=0\pmod M,\\
0,& \Delta\ne0\pmod M.
\end{cases}
\]
Thus the finite grid projects onto $\Delta=0$ only modulo $M$.  In the
order-$\Jpm^2$ and order-$\Jpm^3$ row data, the nonzero components of
$\bDelta=2\bdelta$ are $\pm1$.  Therefore the two-point grid $M=2$ already
distinguishes the finite-size sectors from the clean zero-transport sector.

This is a pure aliasing issue.  Consider one spurious four-loop component with
$\delta_\mu=-\frac12$.  With the effective-row label
$\Delta_\mu=2\delta_\mu=-1$, the two-point average gives
\[
\frac12\left(1+e^{-i\pi(-1)}\right)=0,
\]
so the row is removed.  The perturbative row-table check on the $16$-site
cluster gives exactly this:
\[
\begin{array}{c|c|c}
\text{integer label} & M & \text{relative error to strict }\bdelta=0\\
\hline
2\bdelta & 2 & 0\\
\end{array}
\]
Thus the minimal character projection for this campaign is $2\bdelta$ with
$M=2$.
\end{remark}

\begin{remark}[Why physical twists were insufficient]
A physical flat vector potential gives phases of the form
$e^{-i\bA\cdot\bdelta_r}$.  On the $16$-site cluster the spurious four-loops
have half-cell transported dipoles, so the corresponding exponent is
half-integer in $2\pi$ units.  Physical twist averaging therefore suppresses
these rows but does not project them out exactly.  The auxiliary character
field in Eq.~\eqref{eq:delta_character_phase} is the finite-cluster source
conjugate to the integer transported dipole $\bDelta=2\bdelta$.
\end{remark}

\begin{corollary}[Clean ice-manifold Hamiltonian]
The artifact-free gauge Hamiltonian is
\begin{equation}
\boxed{
\Hclean(\Jpm)=
\sum_{r:\,\bdelta_r=0} c_r\Jpm^{k_r}|t_r\rangle\langle s_r| .
}
\label{eq:hclean}
\end{equation}
It keeps contractible hexagons and removes winding four-loops and wrapping
hexagons.
\end{corollary}

\section{The microscopic spine: structure theorems and the
polarization-parity mask}
\label{sec:qed-spine}

The perturbative derivation above identifies the topological quantum number
that must be removed: nonzero transported dipole.  The production ED protocol
does not stop at the truncated row table.  Its central object is a
\emph{downfolded low-energy operator expressed in the ice basis}, and the
transported-dipole character projection is applied to that operator before
any thermodynamics are computed.  Historically this was implemented as an
eight-point character average of QED low-band extractions; the structure
theorems of Sec.~\ref{sec:structure_theorems} collapse it to
\[
\boxed{
\Hmic
\xrightarrow{\ \mathrm{downfold\ at\ }\btheta=0\ }
\HB
\xrightarrow{\ \mathrm{parity\ mask}\ }
\HB^{\rm clean}
\xrightarrow{\ \mathrm{ED}\ }
C(T),\ S(T),\ \text{gauge-sector observables},
}
\]
a single untwisted computation followed by an exact entrywise projection.
Perturbation theory supplies the character label and the weak-coupling
benchmark; the microscopic solver supplies the dressed low-band operator; the
mask supplies the projection.

\subsection{Gauge-theory meaning}

The pyrochlore lattice is the line graph of the diamond lattice.  A spin on a
pyrochlore site is a directed electric field on a diamond link,
\[
E_\ell=S_i^z,
\]
up to an orientation convention.  The ice rule is Gauss' law,
\begin{equation}
(\nabla\cdot E)_r=\sum_{\ell\supset r}\eta_{r\ell}E_\ell=0.
\label{eq:gauss_law_new}
\end{equation}
so the ice manifold is the zero-charge sector of a compact $U(1)$ gauge
theory.  Integrating out virtual spinons generates local magnetic dynamics,
\begin{equation}
\Heff^{\mathrm{bulk}}
=-g_{\mathrm{hex}}\sum_{p\in\mathrm{hexagons}}
\cos\!\left[(\nabla\times A)_p\right]+\cdots ,
\label{eq:gauge_ring_new}
\end{equation}
where $A_\ell$ is conjugate to $E_\ell$.  This is a local plaquette term: the
physical magnetic ring is a contractible diamond hexagon.

On a three-torus there are also noncontractible Wilson loops, classified by
\[
H_1(T^3,\Z)=\Z^3.
\]
In the thermodynamic gauge theory these are global tunnelling events, not
local photon dynamics.  On a tiny periodic cluster, however, a noncontractible
Wilson loop can be only four spins long.  The protocol removes these
finite-size topological tunnelling events while retaining the local
contractible gauge dynamics.

\begin{remark}[Homotopy fixes the scope of any cleaning scheme]
\label{rem:homotopy_scope}
The torus is aspherical: $T^3=K(\Z^3,1)$, so $\pi_n(T^3)=0$ for $n\ge2$ and
free homotopy classes of closed loops are in bijection with
$\pi_1(T^3)=H_1(T^3,\Z)=\Z^3$.  Two consequences frame everything below.
\emph{(i) Completeness.}  The winding class is the \emph{only} topological
invariant of a completed low-energy process; there is no torsion and no
higher obstruction.  A cleaning scheme keyed to topological labels can remove
at most the nontrivial classes, and the transported-dipole projection removes
exactly these: it is the \emph{maximal} topologically defined cleaning.
\emph{(ii) Matched pairs are honestly trivial.}  A process that winds by
$+\bw$ and later by $-\bw$ lifts to a closed loop in the covering space
$\mathbb R^3$; it has a genuine infinite-lattice counterpart (a virtual
spinon excursion of length $\sim L$) and \emph{cannot} be distinguished from
bulk physics by any superselection label.  Such contributions are ordinary
finite-size image corrections, to be controlled by comparing cluster sizes,
not pathologies to be projected.  This resolves, as a theorem rather than an
observation, why matched-winding contributions survive every averaging
scheme.

At the group level, the transported dipoles generate an abelian group
$\Gamma\subset\tfrac12\Z^3$ sitting in the exact sequence
\[
0\;\to\;\Z^3\bL\ (\text{windings},\ \pi_1)\;\to\;\Gamma\;\to\;
\Gamma/\Z^3\bL\ (\text{polarization classes})\;\to\;0 .
\]
Ordinary boundary twists realize only the characters of the winding subgroup
(flat $U(1)$ connections, $\mathrm{Hom}(\pi_1,U(1))$); the polarization
phases of Sec.~\ref{sec:structure_theorems} realize the characters of the
quotient.  Averaging over a subgroup of characters projects onto the
annihilator of that subgroup --- which is why twist averaging alone keeps
half of every winding channel, and why the gauge-invariant label
$\bdelta=\brho+\bN\bL$ must glue the two factors together.
\end{remark}

\begin{intuition}
The physical photon is built from local magnetic curls around hexagons.  The
artifact is a Wilson loop that wraps around the torus but looks short because
the torus is small.  The transported-dipole character distinguishes these
topological classes even when the spin flips look deceptively similar.
\end{intuition}

\subsection{Character-deformed microscopic Hamiltonian}

For a nearest-neighbor bond, write the minimum-image displacement as
\[
\bd_{ij}=\br_j-\br_i-\bn_{ij}\bL ,
\]
where $\bn_{ij}$ records which periodic image was used.  The microscopic move
$S_i^+S_j^-$ raises site $i$ and lowers site $j$.  The auxiliary
transported-dipole character deformation is
\begin{equation}
\Hmic^{2\delta}(\btheta)=
\Jzz\sum_{\langle ij\rangle}S_i^zS_j^z
-\Jpm\sum_{\langle ij\rangle}
\left[
e^{+i\btheta\cdot(2\bd_{ij})}S_i^+S_j^-
+e^{-i\btheta\cdot(2\bd_{ij})}S_i^-S_j^+
\right].
\label{eq:mic_dipole_twisted_new}
\end{equation}
This is not an ordinary physical boundary flux.  It is a finite-cluster source
field conjugate to the integer transported dipole
\[
\bDelta=2\bdelta .
\]
For the leading finite-size sectors on the cubic-16 and FCC-32 clusters,
$\Delta_\mu=0,\pm1$.  Thus the minimal grid is
\begin{equation}
\theta_\mu\in\{0,\pi\},\qquad M=2,
\label{eq:m2_grid_new}
\end{equation}
with eight character points.  If higher-order or larger-cluster diagnostics
reveal additional transported-dipole labels, $M$ must be increased until the
low-band spectrum and relevant matrix elements converge.

\subsection{Structure theorems: the character average is a mask}
\label{sec:structure_theorems}

The deformation \eqref{eq:mic_dipole_twisted_new} looks like it requires a
separate microscopic diagonalization at every character point.  It does not.
Let
\begin{equation}
\hat{\bm X}=\sum_i \br_i\,S_i^z
\end{equation}
be the many-body polarization in home-cell coordinates, a \emph{diagonal}
operator in the $S^z$ product basis, and define the diagonal unitary
$V(\btheta)=e^{2i\btheta\cdot\hat{\bm X}}$.

\begin{proposition}[Isospectrality at the minimal grid]
\label{prop:isospectral}
Conjugating $\Hmic$ by $V(\btheta)$ multiplies the hop $S_i^+S_j^-$ by
$e^{2i\btheta\cdot(\br_i-\br_j)}$, while the deformation
\eqref{eq:mic_dipole_twisted_new} carries $e^{2i\btheta\cdot\bd_{ij}}$ with
$\bd_{ij}=\br_j-\br_i-\bn_{ij}\bL$; the two differ by the boundary factor
$e^{-2i\btheta\cdot\bn_{ij}\bL}$.  Because $\bL$ is an integer matrix, this
factor equals $1$ at every $\btheta\in\{0,\pi\}^3$.  Hence
\[
\Hmic^{2\delta}(\btheta)=V(\btheta)\,\Hmic\,V(\btheta)^\dagger
\qquad\text{exactly, at all eight corners.}
\]
Consequently all eight corner spectra coincide, and every observable
diagonal in the $S^z$ basis --- including $P_\ice$, all Ising correlators,
and every thermal trace $\Tr f(\Hmic^{2\delta}(\btheta))$ --- is exactly
$\btheta$-independent.
\end{proposition}

Proposition~\ref{prop:isospectral} turns two empirical negative controls
into theorems: observable averaging reproduced the bare peak \emph{to all
digits} because the character family is isospectral, and the pullback
diagnostics were corner-independent because the Gram matrix is.

\begin{proposition}[Mask identity]
\label{prop:mask}
$P_\ice$ commutes with $V(\btheta)$ (both diagonal), so the low-band
ice-frame data of Sec.~\ref{sec:qed-spine} transform covariantly:
\[
X(\btheta)=D(\btheta)X(0),\quad
S(\btheta)=S(0),\quad
\HB(\btheta)=D(\btheta)\,\HB(0)\,D(\btheta)^\dagger,
\]
with $D(\btheta)_{\alpha\alpha}=e^{2i\btheta\cdot\bm x_\alpha}$ and
$\bm x_\alpha=\sum_i\br_i S_i^z(\alpha)$ the polarization of ice state
$\alpha$.  The eight-corner average is therefore an exact entrywise
projection,
\begin{equation}
\boxed{
\left(\HB^{\rm clean}\right)_{\alpha\beta}
=\left(\HB(0)\right)_{\alpha\beta}\,
\Big[\,2(\bm x_\alpha-\bm x_\beta)\ \text{componentwise even}\,\Big],
}
\label{eq:mask_identity}
\end{equation}
the \emph{polarization-parity mask}.  One untwisted downfolding replaces the
eight-point character sweep, and the projection is exact by construction.
\end{proposition}

\begin{proposition}[Factorization of the character group]
\label{prop:factorization}
For general $\btheta$ the same algebra gives
\[
\Hmic^{2\delta}(\btheta)
=V(\btheta)\,H_{\rm twist}\!\left(\bphi=2\bL^{T}\btheta\right)V(\btheta)^\dagger:
\]
every character deformation is gauge-equivalent to an ordinary boundary
twist composed with a polarization rotation, and the full character average
factorizes as
\[
\text{(character average over $\hat\Gamma$)}
=\text{(ordinary twist average)}\circ\text{(parity mask)} .
\]
At $M=2$ the twist factor is trivial and the average is a pure mask; for
$M\ge3$ the twists are genuine and separate diagonalizations are required.
This also explains \emph{structurally} why plain boundary-twist averaging
retains exactly half of every winding channel: it implements only the first
factor of the product.
\end{proposition}

All three statements are verified to machine precision against the stored
eight-corner data by \texttt{notes/verify\_mask\_identity.py}: corner-spectra
spread $3.6\times10^{-15}$; zero mask mismatches on all $8100$ band-operator
entries with residual $<8\times10^{-15}$ at $\Jpm=-0.03$ and $-0.05$
($1284$ of $8100$ entries kept); and, as the negative control, the
dipole-$4\delta$, $M=3$ family --- which contains genuine twists --- has
corner-spectra spread $5\times10^{-2}$ and is \emph{not} reproduced by any
mask.

\begin{intuition}
The character field looks like eight different Hamiltonians, but at the
minimal grid they are all the same Hamiltonian written in eight different
phase conventions.  All of the cleaning power therefore lives in one place:
the fixed ice basis in which the operators are compared.  Rotating each copy
back into that common frame and averaging simply erases every matrix element
whose end states differ by an odd polarization --- which is precisely the
fingerprint the winding four-loop cannot hide.
\end{intuition}

\subsection{Pullback of the QED low band}

Fix one classical ice basis
\[
\{|\alpha\rangle\}_{\alpha=1}^{N_{\ice}},\qquad
P_\ice=\sum_{\alpha=1}^{N_\ice}|\alpha\rangle\langle\alpha|.
\]
For every character point, QED solves
\begin{equation}
\Hmic^{2\delta}(\btheta)|\psi_n(\btheta)\rangle
=E_n(\btheta)|\psi_n(\btheta)\rangle,
\qquad n=1,\ldots,N_\ice ,
\label{eq:qed_low_eigenproblem}
\end{equation}
selecting the lowest ice-like band.  Define
\begin{equation}
X_{\alpha n}(\btheta)=\langle\alpha|\psi_n(\btheta)\rangle .
\label{eq:X_projection_new}
\end{equation}
The projected vectors $P_\ice|\psi_n\rangle$ are not exactly orthonormal.
Their Gram matrix is
\begin{equation}
S(\btheta)=X^\dagger(\btheta)X(\btheta).
\label{eq:gram_matrix_new}
\end{equation}
The first control diagnostic is that $S(\btheta)$ must be nonsingular and well
conditioned.  Physically, this says the selected microscopic band is still the
dressed ice manifold.

Use the Loewdin pullback
\begin{equation}
Q(\btheta)=X(\btheta)S^{-1/2}(\btheta),
\qquad Q^\dagger(\btheta)Q(\btheta)=1.
\label{eq:lowdin_pullback_new}
\end{equation}
The microscopic low-band operator represented in the fixed ice basis is
\begin{equation}
\boxed{
\HB(\btheta)=
Q(\btheta)\,
\mathrm{diag}\!\left(E_1(\btheta),\ldots,E_{N_\ice}(\btheta)\right)
Q^\dagger(\btheta).
}
\label{eq:band_operator_new}
\end{equation}
This is not a finite-order Schrieffer--Wolff truncation.  It contains the
higher-order low-band dressing resolved by the microscopic QED diagonalization
at the chosen coupling.

\begin{remark}[What the Loewdin pullback means]
The word ``pullback'' is being used in a finite-dimensional Hilbert-space
sense.  The microscopic eigenvectors
$|\psi_n(\btheta)\rangle$ live in the full Hilbert space, while the gauge
Hamiltonian should act in the ice basis.  Projecting the microscopic vectors
with $P_\ice$ gives the columns of $X$:
\[
X_{\alpha n}=\langle \alpha|\psi_n\rangle .
\]
These projected vectors are generally not orthonormal.  Their Gram matrix
\[
S_{mn}=\sum_\alpha X^*_{\alpha m}X_{\alpha n}
=\langle\psi_m|P_\ice|\psi_n\rangle
\]
is the metric induced on the projected low-band frame by the ice-manifold
inner product.  Multiplication by $S^{-1/2}$ removes this metric distortion:
\[
Q^\dagger Q
=S^{-1/2}X^\dagger X S^{-1/2}
=1 .
\]
Thus $Q=XS^{-1/2}$ is the orthonormal frame in the ice Hilbert space that is
closest, in the symmetric Loewdin sense, to the projected microscopic low-band
vectors.

This is geometric, but not differential geometry in the sense of curvature or
differential forms.  It is the geometry of finite-dimensional Hilbert-space
subspaces: projection, induced metric, and choice of orthonormal frame.  The
important diagnostic is $\lambda_{\min}S$.  If it is small, at least one
microscopic low-band direction has little ice weight, the pullback amplifies
noise, and the gauge-sector Hamiltonian is not controlled.  If $S$ is well
conditioned, the QED low band is a dressed version of the ice manifold and
Eq.~\eqref{eq:band_operator_new} is a meaningful representation of that dressed
band as an operator on the ice basis.
\end{remark}

\subsection{The clean operator}

The finite-size-clean gauge-sector Hamiltonian is the trivial
transported-dipole character component of Eq.~\eqref{eq:band_operator_new}.
By Proposition~\ref{prop:mask} it is computed \emph{without any character
sweep}:
\begin{equation}
\boxed{
\HB^{\rm clean}
=\mathrm{mask}\!\left[\HB(\btheta=0)\right],
\qquad
\mathrm{mask}:\ \text{keep }(\alpha,\beta)\iff
2(\bm x_\alpha-\bm x_\beta)\in(2\Z)^3,
}
\label{eq:clean_band_operator_new}
\end{equation}
which coincides identically with the eight-corner average
$\frac18\sum_{\btheta\in\{0,\pi\}^3}\HB(\btheta)$.
Only after this operator-level projection do we diagonalize.  This ordering
is essential.  A single transverse hop leaves the ice manifold; the local
hexagon ring appears only after a complete virtual process returns to the ice
manifold.  Filtering microscopic hops directly can distort the process we
want to keep (the direct microscopic average of
Sec.~\ref{sec:qed_validation_new} fails for exactly this reason: winding
lives in the \emph{composition} of hops --- in the $\pi_1$ class of the
completed path --- so no linear operation on $\Hmic$ can see it).  Filtering
the downfolded operator acts on the net transported dipole of the completed
low-energy process.

In the perturbative limit,
\begin{equation}
\HB^{\rm clean}-\mathrm{const.}
=\Hclean+O(\Jpm^4/\Jzz^3),
\label{eq:qed_to_sw_limit_new}
\end{equation}
provided the character grid resolves the transported-dipole sectors present
through third order.  Away from the strict perturbative limit,
Eq.~\eqref{eq:clean_band_operator_new} is a numerical gauge-sector downfolding,
not a full microscopic all-temperature ED trace.

\subsection{Exact implementation checklist (Stage A)}

The production implementation is:
\begin{enumerate}
\item Build the periodic pyrochlore cluster.  Store every nearest-neighbor
bond $(i,j)$, the minimum-image vector $\bd_{ij}$, and the periodic image
integer $\bn_{ij}$.
\item Enumerate tetrahedra and the classical ice basis
$\{|\alpha\rangle\}_{\alpha=1}^{N_\ice}$; record the polarizations
$\bm x_\alpha=\sum_i\br_i S_i^z(\alpha)$ \emph{in the same coordinate frame
as the cluster builder} and precompute the parity mask of
Eq.~\eqref{eq:mask_identity}.
\item Diagonalize the lowest $N_\ice$ states of the \emph{untwisted}
$\Hmic$ using a solver whose eigenvectors are accurately orthonormal.  A
Krylov or block-Lanczos result is accepted only if its vector-Gram
diagnostics pass.
\item Compute $X$, $S=X^\dagger X$, and $Q=XS^{-1/2}$ at $\btheta=0$.  Record
\[
\lambda_{\min}S,\qquad
\lambda_{\max}S,\qquad
\|\Psi^\dagger\Psi-1\|.
\]
\item Assemble $\HB=\HB(\btheta=0)$ from Eq.~\eqref{eq:band_operator_new}
and apply the mask, Eq.~\eqref{eq:clean_band_operator_new}.
\item Diagonalize $\HB^{\rm clean}$ and compute
\[
C(T)=\frac{\langle E^2\rangle_T-\langle E\rangle_T^2}{T^2},
\qquad
S(T)=\ln Z+\frac{\langle E\rangle_T}{T}.
\]
\item Compare against the bare low-band operator
$\HB^{\rm bare}=\HB(\btheta=0)$.  The bare operator should show the winding
four-loop scale, while the clean operator should track $\ghex$ in the
perturbative window.
\item If double-winding (even nonzero transport) contamination is suspected,
refine with genuine boundary twists $\bphi=2\bL^{T}\btheta$ on the $M=4$
grid (Proposition~\ref{prop:factorization}) and check stationarity of the
clean spectrum.
\end{enumerate}

\section{Validation of the QED spine on cubic-16}
\label{sec:qed_validation_new}

The $16$-site cubic cluster is small enough that every piece of the protocol
can be tested against exact microscopic ED.  The fixed-$S^z_{\rm tot}=0$
Hilbert space has dimension
\[
\binom{16}{8}=12870,
\]
and the ice band has $N_\ice=90$ states.

\subsection{Choosing the character}

The perturbative row-table test establishes the minimal character grid.  On
this cluster, physical flat twists and integer winding projectors do not equal
the strict $\bdelta=0$ row projector.  The $2\bdelta$ character does:
\begin{center}
\footnotesize
\begin{tabular}{@{}p{0.30\linewidth}p{0.14\linewidth}p{0.20\linewidth}p{0.22\linewidth}@{}}
\toprule
projector & $T_{\rm peak}/\Jzz$ & rel. error to $\bdelta=0$
& comment\\
\midrule
all rows & $0.0100564$ & $11.6476$ & winding four-loop present\\
integer $N=0$ & $0.0046462$ & $5.8238$ & boundary winding only\\
physical $\bA\!\cdot\!\bdelta$, $M=6$ & $0.0016563$ & $3.0740$
& suppresses, not exact\\
$2\bdelta$ character, $M=2$ & $0.0016045$ & $0$
& minimal row-level charge\\
\bottomrule
\end{tabular}
\end{center}
This table is a row-level proof of the character.  The production calculation
is the QED low-band construction in Eq.~\eqref{eq:clean_band_operator_new}.

\subsection{QED low-band results}

At $\Jpm=-0.05\Jzz$, the historical QED run used all eight $2\bdelta$, $M=2$
character points; those stored operators are now also the machine-precision
verification of Propositions~\ref{prop:isospectral} and~\ref{prop:mask}
(identical corner spectra; averaged operator equal to the mask of the
$\btheta=0$ operator).  The pullback was well conditioned:
\[
\min_{\btheta}\lambda_{\min}S(\btheta)=0.8872896692,
\]
and the saved QED eigenvectors were orthonormal to $2.3\times10^{-14}$.  The
specific-heat peaks are
\begin{center}
\begin{tabular}{lcc}
\toprule
operator & $T_{\rm peak}/\Jzz$ & comment\\
\midrule
QED low band, $\btheta=0$ & $0.0129476$ & bare four-loop scale\\
QED low-band operator average & $0.0012227$ & clean gauge scale\\
perturbative all rows & $0.0100564$ & four-loop artifact\\
perturbative $\bdelta=0$ & $0.0016045$ & leading hexagon target\\
\bottomrule
\end{tabular}
\end{center}

A second full-QED check at $\Jpm=-0.03\Jzz$ confirms that the clean operator
tracks the cubic hexagon scale:
\begin{center}
\begin{tabular}{ccccc}
\toprule
$\Jpm/\Jzz$ & $T_{\rm peak}^{\rm clean}$ & $\ghex$ &
$T_{\rm peak}^{\rm clean}/\ghex$ & $T_{\rm peak}^{\rm bare}$\\
\midrule
$-0.03$ & $0.0002996$ & $0.000324$ & $0.925$ & $0.005541$\\
$-0.05$ & $0.0012227$ & $0.001500$ & $0.815$ & $0.012948$\\
\bottomrule
\end{tabular}
\end{center}
For the $\Jpm=-0.03\Jzz$ run, the minimum ice-subspace overlap was $0.948$
and the full-vector Gram deviation was below $2.1\times10^{-14}$.

\begin{figure}[htbp]
\centering
\includegraphics[width=0.78\textwidth]{fig_qed_dipole2_M2_specific_heat.pdf}
\caption{\textbf{Specific heat from the adopted QED low-band operator
average.}  The orange solid curve is the untwisted microscopic low-band
operator $\HB(\btheta=0)$.  The green solid curve is the operator-level
$2\bdelta$, $M=2$ character average.  Dotted curves are perturbative
ice-manifold matrices.  The QED operator average moves the low-temperature
peak from the spurious four-loop scale toward the clean hexagon scale.}
\label{fig:qed_dipole2_m2_specific_heat_new}
\end{figure}

\subsection{Rejected shortcuts}

The naive direct microscopic average
\[
\Hmic^{\rm avg}=\frac18\sum_{\btheta}\Hmic^{2\delta}(\btheta)
\]
was tested separately.  It produces a low band with large ice overlap
($0.982$ mean at $\Jpm=-0.05\Jzz$), but its specific-heat peak follows
$\gfour$, not $\ghex$:
\begin{center}
\begin{tabular}{ccc}
\toprule
$\Jpm/\Jzz$ & $T_{\rm peak}/\gfour$ & $T_{\rm peak}/\ghex$\\
\midrule
$-0.02$ & $0.125$ & $2.08$\\
$-0.03$ & $0.119$ & $1.32$\\
$-0.04$ & $0.114$ & $0.946$\\
$-0.05$ & $0.109$ & $0.725$\\
$-0.06$ & $0.104$ & $0.581$\\
\bottomrule
\end{tabular}
\end{center}
It changes the coefficient of the finite-size artifact but does not isolate
the hexagon manifold.

Observable averaging also fails.  At $\Jpm=-0.05\Jzz$, averaging the
twist-resolved low-band specific heats gives
\[
T_{\rm peak}/\Jzz=0.0129476,
\]
the same peak as the bare periodic low band.  These shortcuts are useful
negative controls, but the spine of the method is
Eq.~\eqref{eq:clean_band_operator_new}.

\section{Graded Brillouin--Wigner downfolding: the protocol at any $\Jpm$}
\label{sec:bw_protocol}

The Loewdin construction of Sec.~\ref{sec:qed-spine} anchors the lowest
$N_\ice$ \emph{eigenstates} to the ice basis, and therefore requires the
microscopic low band to still \emph{be} an ice band.  That requirement fails
at material-scale couplings (the pullback Gram matrix becomes singular by
$\Jpm/\Jzz\simeq-0.18$ on cubic-16).  The all-$\Jpm$ engine replaces
eigenstate selection by the exact Feshbach--Schur (Brillouin--Wigner) map
onto the ice block: with $P=P_\ice$, $Q=1-P$,
\begin{equation}
\boxed{
F(z)=PHP+PHQ\,(z-QHQ)^{-1}\,QHP ,
}
\label{eq:bw_map}
\end{equation}
defined for every $\Jpm$ and every $z\notin\mathrm{spec}(QHQ)$.  No
eigenvectors, no band selection, no Gram matrix.  The geometric expansion of
the resolvent is a sum over \emph{completed} virtual excursions
$P\to Q\to\cdots\to Q\to P$, each carrying a definite transported dipole, so
the character grading of Sec.~\ref{sec:transported-dipole} applies to $F(z)$
term by term, resummed to infinite order.  Eigenvalues of $H$ with nonzero
ice weight are exactly the solutions of the nonlinear problem
$\det[F(E)-E]=0$.

\subsection{Equivariance and the masked operator}

$P$ and $V(\btheta)$ are both diagonal in the $S^z$ basis, so the structure
theorems transfer verbatim to the BW map:
\[
F(z;\btheta)=D(\btheta)\,F_{\rm twist}\!\left(z;2\bL^{T}\btheta\right)
D(\btheta)^\dagger ,
\]
and at the minimal grid the clean object is a single masked operator,
\begin{equation}
\boxed{
F_{\rm clean}(z)=\mathrm{mask}\!\left[F(z;\btheta=0)\right].
}
\label{eq:bw_masked}
\end{equation}
This projects every virtual process, at every order in $\Jpm$, onto
$\bDelta=2\bdelta=0\pmod 2$: the winding four-loops and wrapping hexagons are
removed identically, at any coupling, whether or not perturbation theory
converges and whether or not the low band is ice-like.

\subsection{Self-consistent solution and built-in diagnostics}

The clean gauge-sector levels solve
\begin{equation}
E_n=\mathrm{eig}_n\,F_{\rm clean}(E_n),
\label{eq:bw_selfconsistent}
\end{equation}
found by a damped fixed-point iteration per level.  Three diagnostics come
for free:
\begin{itemize}
\item \textbf{Non-ice weight.}  At a solution with ice-block eigenvector
$u_n$, the full eigenstate carries non-ice weight
\[
w_n^{\rm out}
=\frac{u_n^\dagger\,PHQ\,(E_n-QHQ)^{-2}\,QHP\,u_n}
       {1+u_n^\dagger\,PHQ\,(E_n-QHQ)^{-2}\,QHP\,u_n}
\quad\in[0,1),
\]
which equals the $z$-derivative of the eigenvalue of $F$ and is therefore
also the convergence rate of the fixed point: the solver converges exactly
as fast as the state is ice-like.
\item \textbf{Pole distance.}  $\min_k|E_n-\lambda_k(QHQ)|$ measures the
spectral isolation of the gauge-sector resonance from the spinon continuum.
\item \textbf{Anchoring check.}  The \emph{unmasked} self-consistent BW
spectrum must reproduce the exact lowest $N_\ice$ microscopic levels
whenever the low band is the ice band; a mismatch is the sharp statement
that eigenstate-anchored constructions (Stage A) have lost their target,
while the BW resonances remain defined.
\end{itemize}
Breakdown is therefore \emph{interpretable}: when the solutions collide with
$\mathrm{spec}(QHQ)$ and the non-ice weight approaches one, the statement is
that the cluster's gauge sector has dissolved into the spinon continuum at
that coupling --- a physical finding, not a numerical failure.

\section{The cubic-16 validation ladder}
\label{sec:jpm_validity_classification_new}

Every ingredient can be tested exactly on cubic-16
(\texttt{notes/masked\_bw\_16site.py}; records in
\texttt{notes/masked\_bw\_16site\_results.json}).  Per coupling the script
computes the full microscopic spectrum densely (ground truth), the Stage-A
masked Loewdin operator where it exists, and the Stage-B masked BW solution
with all diagnostics.  Convention anchors: the exact lowest $90$ levels agree
with the stored eight-corner files to $3\times10^{-14}$, and the
$\btheta=0$ Gram eigenvalue $\eta_{\min}$ reproduces the earlier validity
sweep at every coupling.

\begin{center}
\footnotesize
\begin{tabular}{@{}rrrrrrrl@{}}
\toprule
$\Jpm/\Jzz$ & $\eta_{\min}$ & conv. & $T^{\rm clean}_{\rm peak}/\ghex$ &
$\max_n w_n^{\rm out}$ & pole gap & Loewdin $/\ghex$ & verdict\\
\midrule
$-0.03$ & $0.948$ & $90/90$ & $0.927$ & $0.029$ & $0.84$ & $0.927$ &
controlled\\
$-0.05$ & $0.887$ & $90/90$ & $0.859$ & $0.074$ & $0.74$ & $0.817$ &
controlled\\
$-0.08$ & $0.790$ & $90/90$ & $0.769$ & $0.194$ & $0.59$ & $0.648$ &
gray zone\\
$-0.10$ & $0.724$ & $90/90$ & $0.705$ & $0.241$ & $0.50$ & $0.540$ &
gray zone\\
$-0.18$ & $0.360$ & $90/90$ & $0.435$ & $0.536$ & $0.185$ & $(0.257)$ &
recovered; strongly dressed\\
$-0.30$ & $\sim0$ & $52/90$ & --- & $0.9995$ & $2.1\times10^{-4}$ & --- &
gauge sector dissolved\\
\bottomrule
\end{tabular}
\end{center}

Three conclusions, in increasing order of importance:

\begin{enumerate}
\item \textbf{Reduction.}  In the controlled window the masked BW solution
coincides with the eight-corner character average.  At $\Jpm=-0.03$ the two
clean specific-heat peaks are identical to all reported digits; the residual
maximum spectral deviation ($5\times10^{-4}$) is a uniform shift of the whole
band, invisible to thermodynamics.  At $-0.05$ the peaks agree to $5\%$, and
the spread between the two constructions tracks the non-ice weight --- the
diagnostic doubles as the error bar between resummations.
\item \textbf{Extension.}  The CHO-like point $\Jpm/\Jzz=-0.18$, where the
Gram matrix collapse ($\eta_{\min}=0.36$) invalidated Stage A, is
\emph{recovered}.  The anchoring check makes the mechanism explicit: the
unmasked BW resonances differ from the exact lowest-$90$ spectrum by up to
$4.2\,\Jzz$ there (five levels do not even converge to the lowest band), so
the lowest eigenstates are simply no longer the ice band; but the ice-block
resonances themselves still exist, all $90$ masked levels converge, and they
remain separated from the spinon continuum by $0.185\,\Jzz$.  The old
``failure'' was an artifact of eigenstate anchoring, not of the physics.
\item \textbf{Honest boundary.}  At the CZO-like point $\Jpm/\Jzz=-0.30$ the
diagnostics deliver a verdict instead of a crash: $38$ of $90$ levels do not
converge, the surviving solutions sit $2\times10^{-4}\,\Jzz$ from the
continuum poles, and the non-ice weight reaches $0.9995$.  On this cluster,
at this coupling, there is no isolated gauge sector to clean.  Whether
FCC-32 restores one is now a well-posed question for the same diagnostics.
\end{enumerate}

\begin{figure}[htbp]
\centering
\includegraphics[width=0.62\textwidth]{fig_bw_peaks.pdf}
\caption{\textbf{Peak scaling across the full coupling range (cubic-16).}
The bare low band (orange) tracks the winding artifact $\gfour$; the masked
BW clean operator (green) tracks the photon scale $\ghex$; crosses are the
masked Loewdin operator where its Gram matrix is nonsingular.}
\label{fig:bw_peaks}
\end{figure}

\begin{figure}[htbp]
\centering
\includegraphics[width=0.95\textwidth]{fig_bw_ct.pdf}
\caption{\textbf{Clean specific heat in and beyond the old validity window.}
Left: the controlled point $\Jpm=-0.05\Jzz$.  Right: the CHO-like point
$\Jpm=-0.18\Jzz$, inaccessible to the eigenstate-anchored protocol, where
the masked BW operator still yields fully converged clean thermodynamics
(reported with its $54\%$ maximum dressing diagnostic).  Dashed and dotted
verticals mark $\gfour$ and $\ghex$.}
\label{fig:bw_ct}
\end{figure}

\begin{figure}[htbp]
\centering
\includegraphics[width=0.95\textwidth]{fig_bw_diagnostics.pdf}
\caption{\textbf{All-$\Jpm$ validity diagnostics.}  Left: the old Stage-A
gate $\eta_{\min}$ (purple) against the BW non-ice weight (green); the
former collapses at $-0.18$ while the latter degrades smoothly and only
saturates at $-0.30$.  Right: minimum distance of the solved levels to
$\mathrm{spec}(QHQ)$ --- the spectral-isolation diagnostic whose collapse at
$-0.30$ constitutes the physical statement that the gauge sector has merged
with the spinon continuum.}
\label{fig:bw_diagnostics}
\end{figure}

\begin{figure}[htbp]
\centering
\includegraphics[width=0.82\textwidth]{fig_jpm_validity_sweep_16site_summary.pdf}
\caption{\textbf{Historical Stage-A validity sweep} (Loewdin pullback
$\eta_{\min}$ and eight-corner clean peaks), retained for comparison with
the BW ladder above.}
\label{fig:jpm_validity_sweep_16site_new}
\end{figure}

The status classification is now three-tiered and applies at any coupling:
\emph{controlled} ($w^{\rm out}\lesssim0.1$, pole gap $O(\Jzz)$: production
gauge-sector ED), \emph{gray zone / exploratory}
($0.1\lesssim w^{\rm out}\lesssim0.5$: report with the BW--Loewdin spread as
the method error), and \emph{dissolved} (unconverged levels, pole gap
$\to0$, $w^{\rm out}\to1$: the finding is the absence of an isolated gauge
sector, and only a larger cluster can change it).  The sign of $\Jpm$ plays
no role in the classification; the zero-flux side additionally admits
sign-free QMC benchmarks.

\section{All-temperature consistency check on cubic-16}
\label{sec:all_temperature_check_new}

The QED low-band operator average is a gauge-sector construction.  By itself
it does not contain the larger-temperature peak produced by sectors with
violated tetrahedron constraints.  Therefore the all-temperature diagnostic
must alter only the part of the spectrum where the finite-size pathology lives.

Let
\[
E^{\rm full}_1\le E^{\rm full}_2\le \cdots \le E^{\rm full}_{2^{16}}
\]
be the exact untwisted microscopic spectrum, including all $S^z_{\rm tot}$
sectors, and let
\[
\epsilon^{\rm clean}_1\le\cdots\le \epsilon^{\rm clean}_{N_{\rm ice}}
\]
be the spectrum of $\HB^{\rm clean}$.  Define
\begin{equation}
\label{eq:allT_spectral_replacement_new}
\widetilde E_a =
\begin{cases}
\epsilon^{\rm clean}_a+s,&1\le a\le N_{\rm ice},\\
E^{\rm full}_a,&a>N_{\rm ice},
\end{cases}
\qquad
s =
\frac{1}{N_{\rm ice}}\sum_{a=1}^{N_{\rm ice}}E^{\rm full}_a
-
\frac{1}{N_{\rm ice}}\sum_{a=1}^{N_{\rm ice}}\epsilon^{\rm clean}_a .
\end{equation}
The shift $s$ preserves the center of the ice band relative to the spinon
continuum and changes only the internal gauge-sector splittings.

QED exact full-spectrum ED gives all $2^{16}=65536$ eigenvalues.  The lowest
$90$ exact eigenvalues agree with the untwisted low-band extraction to
$6.22\times10^{-14}$.  The peak comparison is
\begin{center}
\begin{tabular}{lcc}
\toprule
curve & low-$T$ peak & larger-$T$ peak\\
\midrule
exact full ED & $0.0129475$ & $0.252820$\\
full ED with clean ice-band replacement & $0.0012227$ & $0.252931$\\
\bottomrule
\end{tabular}
\end{center}
The larger peak shifts by only $4.37\times10^{-4}$ relatively, with height
change $-4.83\times10^{-3}$.  The finite-size correction acts where intended:
on the low-energy ice-band winding artifact, not on the microscopic
constraint-violation scale.

\begin{figure}[htbp]
\centering
\includegraphics[width=0.82\textwidth]{fig_full_ed_all_temperature_specific_heat.pdf}
\caption{\textbf{All-temperature full-ED consistency check.}  The black curve
is the exact untwisted microscopic full-ED specific heat for all $65536$ states
of the $16$-site cluster.  The green curve keeps every non-ice eigenvalue fixed
and replaces only the lowest $90$ ice-band eigenvalues by the clean QED
low-band operator spectrum, with the band-center shift in
Eq.~\eqref{eq:allT_spectral_replacement_new}.  The correction moves the
spurious low-temperature four-loop peak while leaving the larger full-ED peak
at $T/\Jzz\simeq0.253$ in place.}
\label{fig:full_ed_all_temperature_specific_heat_new}
\end{figure}

\section{FCC-32 production protocol}
\label{sec:fcc32_protocol_new}

The $2\times2\times2$ FCC pyrochlore cluster has $32$ spins and
\[
N_\ice=2970
\]
ice states.  Full microscopic ED of all $2^{32}$ states is unavailable, so
the production calculation is the gauge-sector construction above.  By
Proposition~\ref{prop:mask} it requires a \emph{single} untwisted low-band
extraction:
\[
H_{\mathrm{band,FCC32}}^{\rm clean}
=\mathrm{mask}\!\left[H_{\mathrm{band,FCC32}}(\btheta=0)\right],
\qquad
H_{\mathrm{band,FCC32}}^{\rm bare}=H_{\mathrm{band,FCC32}}(\btheta=0),
\]
an eightfold cost reduction relative to the character sweep, with the
projection exact rather than grid-sampled.  At material-scale couplings the
Stage-B object replaces the band extraction: $2970$ shifted block-Krylov
solves $(z-QHQ)^{-1}QHP\,e_\alpha$ in the fixed-$S^z$ sector (dimension
$601{,}080{,}390$), one Krylov space per right-hand side serving all $z$,
plus a Lanczos estimate of the $QHQ$ edge for the pole-gap diagnostic.

The FCC implementation must:
\begin{enumerate}
\item express $\bd_{ij}$, $\bDelta=2\bdelta$, and the polarizations
$\bm x_\alpha$ in the FCC primitive coordinate frame used by the cluster
builder, and unit-test the parity mask against the cubic-16 reference
(\texttt{notes/verify\_mask\_identity.py} logic);
\item run the untwisted low-band extraction (Stage A) or the BW solves
(Stage B) with ice-weight, Gram, pole-gap, and non-ice-weight diagnostics;
\item diagonalize both $H_{\mathrm{band,FCC32}}^{\rm bare}$ and
$H_{\mathrm{band,FCC32}}^{\rm clean}$;
\item check that the clean peak scales with $\ghex$ at small $|\Jpm/\Jzz|$;
\item spot-check $M$-convergence with genuine boundary twists
$\bphi=2\bL^{T}\btheta$ on the $M=4$ grid
(Proposition~\ref{prop:factorization});
\item benchmark the zero-flux side against sign-free QMC before making
pi-flux material claims.
\end{enumerate}

The material-facing runs should be interpreted as gauge-sector predictions.
For the first U(1)-symmetric campaign one sets $J_{\pm\pm}=0$ after rotating
away $J_{xz}$ and uses representative pi-flux points,
\[
\text{CZO-like: } \Jzz=1,\quad \Jpm=-0.30,\quad J_{\pm\pm}=0,
\]
\[
\text{CHO-like: } \Jzz=1,\quad \Jpm=-0.18,\quad J_{\pm\pm}=0.
\]
These are not zero-flux material partners.  Zero-flux calculations are generic
sign-free QMC benchmarks for the method.

\end{document}
